\documentclass[UTF8,zihao=-4]{ctexart}

% =========================================================
% Jerry 复习笔记 LaTeX 模板
% 用途：把任意课程的知识点复习讲义整理成正式 PDF
% 编译建议：xelatex -> xelatex
% 默认水印：Jerry 制作｜仅供个人复习使用｜禁止盗印与商用传播
% =========================================================

% =========================
% 页面与基础宏包
% =========================
\usepackage[a4paper,margin=2.15cm]{geometry}
\usepackage{amsmath,amssymb,bm,mathtools}
\usepackage{esint}
\usepackage{booktabs,longtable,tabularx,array,multicol}
\usepackage[most]{tcolorbox}
\usepackage{xcolor}
\usepackage{enumitem}
\usepackage{hyperref}
\usepackage{fancyhdr}
\usepackage{titlesec}
\usepackage{lastpage}
\usepackage{siunitx}
\usepackage{fvextra}
\usepackage{tikz}
\usepackage{eso-pic}

% =========================
% 可修改的文档变量
% =========================
\newcommand{\CourseName}{课程名称}
\newcommand{\NoteTitle}{\CourseName 期末复习笔记}
\newcommand{\NoteSubtitle}{按样卷、例题与复习资料重新整理}
\newcommand{\MainLineText}{模块一 -- 模块二 -- 模块三 -- 综合应用}
\newcommand{\AuthorName}{Jerry}
\newcommand{\Version}{V1.0}
\newcommand{\WatermarkText}{Jerry 制作｜仅供个人复习使用｜禁止盗印与商用传播}

% =========================
% 字体设置
% 若本机没有 Noto 字体，可替换为系统可用中文字体
% =========================
\setCJKmainfont{Noto Serif CJK SC}
\setCJKsansfont{Noto Sans CJK SC}
\setmainfont{DejaVu Serif}
\setsansfont{DejaVu Sans}
\setmonofont{DejaVu Sans Mono}

% =========================
% PDF 元信息与超链接
% =========================
\hypersetup{
  colorlinks=true,
  linkcolor=blue!55!black,
  urlcolor=blue!55!black,
  citecolor=blue!55!black,
  pdftitle={\NoteTitle},
  pdfauthor={Jerry}
}

% =========================
% 主题颜色
% =========================
\definecolor{mainblue}{HTML}{1F4E79}
\definecolor{lightblue}{HTML}{EAF3FF}
\definecolor{warnred}{HTML}{B3261E}
\definecolor{lightred}{HTML}{FFF1F1}
\definecolor{darkgreen}{HTML}{1B6B3A}
\definecolor{lightgreen}{HTML}{EEF9F0}
\definecolor{orange}{HTML}{B26A00}
\definecolor{lightorange}{HTML}{FFF7E8}
\definecolor{softgray}{HTML}{777777}

% =========================
% 章节标题样式
% =========================
\ctexset{
  section={format=\Large\bfseries\sffamily\color{mainblue}},
  subsection={format=\large\bfseries\sffamily\color{mainblue!85!black}},
  subsubsection={format=\normalsize\bfseries\sffamily\color{mainblue!75!black}}
}
\titleformat{\paragraph}[runin]{\bfseries\sffamily\color{mainblue!75!black}}{}{0pt}{}[\quad]

% =========================
% 页眉页脚
% =========================
\pagestyle{fancy}
\fancyhf{}
\lhead{\NoteTitle}
\rhead{\leftmark}
\cfoot{\thepage/\pageref{LastPage}}
\setlength{\headheight}{15pt}

% =========================
% 列表与表格
% =========================
\setlist[itemize]{leftmargin=2em,itemsep=0.25em,topsep=0.25em}
\setlist[enumerate]{leftmargin=2em,itemsep=0.25em,topsep=0.25em}
\renewcommand{\arraystretch}{1.35}
\allowdisplaybreaks

% =========================
% 常用数学命令：按课程需要保留或扩展
% =========================
\newcommand{\dd}{\,\mathrm d}
\newcommand{\ii}{\mathrm i}
\newcommand{\ee}{\mathrm e}
\newcommand{\ve}[1]{\bm{#1}}
\newcommand{\grad}{\nabla}

% 物理类课程常用命令，可按需删除
\newcommand{\eps}{\varepsilon}
\newcommand{\epsz}{\varepsilon_0}
\newcommand{\muz}{\mu_0}
\newcommand{\mur}{\mu_r}
\newcommand{\E}{\ve E}
\newcommand{\B}{\ve B}
\newcommand{\D}{\ve D}
\newcommand{\Hf}{\ve H}
\newcommand{\J}{\ve j}
\newcommand{\dl}{\dd\ve l}
\newcommand{\dS}{\dd\ve S}
\newcommand{\dV}{\dd V}

% =========================
% 彩色盒子环境
% =========================
\newtcolorbox{keybox}[1][]{
  breakable,
  colback=lightblue,
  colframe=mainblue,
  arc=2mm,
  title=\textbf{#1},
  fonttitle=\bfseries\sffamily,
  coltitle=white
}

\newtcolorbox{warnbox}[1][]{
  breakable,
  colback=lightred,
  colframe=warnred,
  arc=2mm,
  title=\textbf{#1},
  fonttitle=\bfseries\sffamily,
  coltitle=white
}

\newtcolorbox{examplebox}[1][]{
  breakable,
  colback=lightgreen,
  colframe=darkgreen,
  arc=2mm,
  title=\textbf{#1},
  fonttitle=\bfseries\sffamily,
  coltitle=white
}

\newtcolorbox{methodbox}[1][]{
  breakable,
  colback=lightorange,
  colframe=orange,
  arc=2mm,
  title=\textbf{#1},
  fonttitle=\bfseries\sffamily,
  coltitle=white
}

% 代码块：程序设计类课程可用
\DefineVerbatimEnvironment{Code}{Verbatim}{fontsize=\small,breaklines,breakanywhere,frame=single,framesep=2mm}

% =========================
% 每页斜向浅灰水印
% 若觉得太密，可调 step 数值、opacity、scale
% =========================
\AddToShipoutPictureBG{%
  \begin{tikzpicture}[remember picture,overlay]
    \foreach \x in {-8,-2,4,10,16,22,28}{%
      \foreach \y in {-4,2,8,14,20,26,32}{%
        \node[rotate=35,scale=1.05,text=softgray,opacity=0.13] at ([xshift=\x cm,yshift=\y cm]current page.south west) {\WatermarkText};
      }
    }
  \end{tikzpicture}%
}

\begin{document}

% =========================
% 封面页
% =========================
\begin{titlepage}
\centering
\vspace*{2.5cm}

{\Huge\bfseries\sffamily \NoteTitle\par}

\vspace{0.8cm}

{\Large \NoteSubtitle\par}

\vspace{0.5cm}

{\large \MainLineText\par}

\vfill

\begin{tcolorbox}[
  colback=lightblue,
  colframe=mainblue,
  width=0.86\textwidth,
  arc=2mm
]
\textbf{使用定位：}
这份笔记不是教材逐章摘抄，而是把样卷、例题和复习资料按期末复习需要重新组织。重点放在核心公式、适用条件、典型题型、解题模板和易错判断。\par

\textbf{复习建议：}
先看“课程主线”和“公式速查”，再按“方法模板”刷题，最后用“易错点汇总”和“速记清单”查漏补缺。
\end{tcolorbox}

\vfill

{\large 资料整理：\AuthorName\par}
{\large 版本：\Version\par}
{\large \today\par}

\end{titlepage}

% =========================
% 目录
% =========================
\tableofcontents
\newpage

% =========================================================
% 第一章：总览
% =========================================================
\section{总览：本课程的主线}

本课程可以按以下逻辑理解：
\[
\text{基础概念} \to \text{核心公式} \to \text{典型模型} \to \text{高频题型} \to \text{综合应用}
\]

\begin{keybox}[本课程最终主线]
\begin{itemize}
  \item 主线一：填写本课程从基础到综合的第一条逻辑链。
  \item 主线二：填写本课程最重要的公式或方法链。
  \item 主线三：填写本课程高频大题的共同套路。
\end{itemize}
\end{keybox}

\begin{warnbox}[最容易混淆的概念]
\begin{enumerate}
  \item 概念 A 与概念 B 的区别：填写区别。
  \item 公式 A 与公式 B 的适用条件：填写条件。
  \item 出现某关键词时不能直接套公式，必须先检查边界、方向、定义域或单位。
\end{enumerate}
\end{warnbox}

% =========================================================
% 第二章：考试结构与复习优先级
% =========================================================
\section{考试结构与复习优先级}

\subsection{题型与分值结构}

\begin{longtable}{p{0.18\textwidth}p{0.14\textwidth}p{0.16\textwidth}p{0.14\textwidth}p{0.28\textwidth}}
\toprule
\textbf{题型} & \textbf{题量} & \textbf{单题分值} & \textbf{总分} & \textbf{考察特点} \\
\midrule
选择题 & 待填 & 待填 & 待填 & 概念判断、快速计算 \\
填空题 & 待填 & 待填 & 待填 & 公式应用、短计算 \\
计算题 & 待填 & 待填 & 待填 & 固定套路、过程得分 \\
解答题 & 待填 & 待填 & 待填 & 综合计算、方法选择 \\
应用题 & 待填 & 待填 & 待填 & 建模、综合应用 \\
\bottomrule
\end{longtable}

\subsection{考点优先级}

\begin{longtable}{p{0.18\textwidth}p{0.22\textwidth}p{0.24\textwidth}p{0.26\textwidth}}
\toprule
\textbf{优先级} & \textbf{模块} & \textbf{常见题型} & \textbf{复习要求} \\
\midrule
第一优先级 & 必考核心 & 大题/计算题 & 必须掌握公式、步骤和易错点 \\
第二优先级 & 高频重点 & 选择/填空/大题 & 熟悉常见变形 \\
第三优先级 & 可能考点 & 小题/应用题 & 会基本判断和套用 \\
第四优先级 & 了解即可 & 概念题 & 不深挖 \\
\bottomrule
\end{longtable}

% =========================================================
% 第三章起：模块化知识点
% 复制本章结构即可新增模块
% =========================================================
\section{模块一：模块名称}

\subsection{本模块考试地位}

本模块属于：\textbf{必考核心 / 高频重点 / 可能考点 / 了解即可}。

常见题型：
\begin{itemize}
  \item 选择题：填写常见考法。
  \item 填空题：填写常见考法。
  \item 大题：填写常见考法。
\end{itemize}

\begin{keybox}[本模块复习目标]
\begin{enumerate}
  \item 会判断：填写判断目标。
  \item 会计算：填写计算目标。
  \item 会区分：填写易混点。
  \item 会应用：填写应用场景。
\end{enumerate}
\end{keybox}

\subsection{核心概念}

用考试语言解释核心概念，不要照抄教材。

\begin{itemize}
  \item \textbf{概念一：}填写考前理解版解释。
  \item \textbf{概念二：}填写考前理解版解释。
  \item \textbf{概念三：}填写考前理解版解释。
\end{itemize}

\subsection{核心公式与结论}

\begin{keybox}[核心公式]
\begin{align*}
\text{公式一：}\quad & A = B + C, \\
\text{公式二：}\quad & F(x)=\int_a^b f(x,t)\dd t, \\
\text{公式三：}\quad & \grad f = \left(\frac{\partial f}{\partial x},\frac{\partial f}{\partial y},\frac{\partial f}{\partial z}\right).
\end{align*}

\textbf{适用条件：}
\begin{itemize}
  \item 条件一：填写公式适用条件。
  \item 条件二：填写不能使用的情况。
  \item 条件三：填写考试常见变形。
\end{itemize}
\end{keybox}

\subsection{常见模型或常用结论}

\begin{longtable}{p{0.28\textwidth}p{0.62\textwidth}}
\toprule
\textbf{模型/题型} & \textbf{结论/做法} \\
\midrule
模型一 & 填写结论或固定处理方式。 \\
模型二 & 填写结论或固定处理方式。 \\
模型三 & 填写结论或固定处理方式。 \\
模型四 & 填写结论或固定处理方式。 \\
\bottomrule
\end{longtable}

\begin{warnbox}[结论公式必须看适用条件]
\begin{itemize}
  \item 公式一只适用于：填写条件。
  \item 公式二不能用于：填写禁用场景。
  \item 如果题目出现：填写关键词，应优先考虑：填写方法。
\end{itemize}
\end{warnbox}

\subsection{方法模板}

\begin{methodbox}[方法模板：题型名称]
\begin{enumerate}
  \item 第一步：识别题型，判断考的是哪个知识点。
  \item 第二步：写出核心公式或核心判别条件。
  \item 第三步：代入题目条件，统一变量、区域、方向或单位。
  \item 第四步：计算并化简。
  \item 第五步：检查方向、符号、单位、定义域、边界或特殊点。
\end{enumerate}
\end{methodbox}

\subsection{典型例题}

\begin{examplebox}[典型例题]
\textbf{题目：}
填写一道与样卷同类但不照搬的例题。

\textbf{思路：}
说明这道题的核心识别点和解题路线。

\textbf{解：}
\[
\text{在这里写主要计算过程。}
\]

\textbf{答案：}
填写最终答案。

\textbf{提醒：}
本题最容易错在：填写易错点。
\end{examplebox}

\subsection{本模块易错点}

\begin{warnbox}[本模块易错点]
\begin{enumerate}
  \item 易错点一：填写错误表现与正确做法。
  \item 易错点二：填写错误表现与正确做法。
  \item 易错点三：填写错误表现与正确做法。
  \item 易错点四：填写错误表现与正确做法。
\end{enumerate}
\end{warnbox}

\subsection{本模块速记}

\begin{keybox}[本模块速记]
\begin{itemize}
  \item 看到……先判断……
  \item 出现……优先使用……
  \item 计算……一定注意……
  \item 最后答案检查……
\end{itemize}
\end{keybox}

% =========================================================
% 高频题型模板
% =========================================================
\section{高频题型方法模板}

\begin{methodbox}[题型一：填写题型名称]
\textbf{识别特征：}
看到题目中出现……，一般就是考……

\textbf{固定步骤：}
\begin{enumerate}
  \item ……
  \item ……
  \item ……
\end{enumerate}

\textbf{易错点：}
\begin{itemize}
  \item ……
  \item ……
\end{itemize}
\end{methodbox}

\begin{methodbox}[题型二：填写题型名称]
\textbf{识别特征：}
……

\textbf{固定步骤：}
\begin{enumerate}
  \item ……
  \item ……
  \item ……
\end{enumerate}
\end{methodbox}

% =========================================================
% 期末高频易错点汇总
% =========================================================
\section{期末高频易错点汇总}

\begin{warnbox}[一页核心易错点]
\begin{enumerate}
  \item 易错点一：……
  \item 易错点二：……
  \item 易错点三：……
  \item 易错点四：……
  \item 易错点五：……
\end{enumerate}
\end{warnbox}

% =========================================================
% 公式速查表
% =========================================================
\section{公式速查表}

\subsection{模块一公式}

\begin{longtable}{p{0.32\textwidth}p{0.58\textwidth}}
\toprule
\textbf{名称} & \textbf{公式 / 结论} \\
\midrule
公式一 & $A=B+C$ \\
公式二 & $F(x)=\int_a^b f(x,t)\dd t$ \\
公式三 & $\grad f=(f_x,f_y,f_z)$ \\
\bottomrule
\end{longtable}

\subsection{模块二公式}

\begin{longtable}{p{0.32\textwidth}p{0.58\textwidth}}
\toprule
\textbf{名称} & \textbf{公式 / 结论} \\
\midrule
公式一 & 填写公式。 \\
公式二 & 填写公式。 \\
公式三 & 填写公式。 \\
\bottomrule
\end{longtable}

% =========================================================
% 考前复习计划与速记
% =========================================================
\section{考前复习计划}

\subsection{七天版}
\begin{methodbox}[7 天复习安排]
\begin{enumerate}
  \item Day 1：过课程主线和核心公式。
  \item Day 2：刷选择题和填空题，整理小题易错点。
  \item Day 3：专项突破计算题。
  \item Day 4：专项突破解答题。
  \item Day 5：专项突破应用题或综合题。
  \item Day 6：完成模拟卷一并订正。
  \item Day 7：完成模拟卷二，背速记清单。
\end{enumerate}
\end{methodbox}

\subsection{三天版}
\begin{methodbox}[3 天冲刺安排]
\begin{enumerate}
  \item Day 1：背核心公式，过必考题型模板。
  \item Day 2：做一套模拟卷，整理错题。
  \item Day 3：复盘错题，背最后速记清单。
\end{enumerate}
\end{methodbox}

\section{考前最后 30 分钟速记清单}

\begin{keybox}[最后 30 分钟只看这些]
\begin{enumerate}
  \item 必背公式一：……
  \item 必背公式二：……
  \item 必背判别条件：……
  \item 大题步骤模板：……
  \item 最容易丢分点：……
\end{enumerate}
\end{keybox}

% =========================================================
% 附录
% =========================================================
\section{附录：补充结论或拓展题型}

\subsection{补充结论一}
填写补充结论。

\subsection{补充结论二}
填写补充结论。

\end{document}
